\chapter{Glucose (Urine)}
\section*{What Is This Test?}
Urine glucose (glucosuria) testing detects the presence of glucose in the urine using a dipstick reagent pad impregnated with glucose oxidase, peroxidase, and a chromogen (potassium iodide or tetramethylbenzidine). The glucose oxidase specifically oxidizes glucose to gluconic acid and hydrogen peroxide; peroxidase catalyzes the oxidation of the chromogen by hydrogen peroxide, producing a color change proportional to the glucose concentration. Normally, glucose is freely filtered by the glomerulus and is completely reabsorbed by the proximal tubular sodium-glucose cotransporters (SGLT2, responsible for 90\% of glucose reabsorption, and SGLT1, 10\%) when the blood glucose is below the renal threshold of approximately 180 mg/dL (10 mmol/L) \cite{moe2020sodium}. Glucosuria occurs when the blood glucose exceeds the renal threshold (hyperglycemic glucosuria, most commonly in diabetes mellitus) or when the renal threshold is normal but the tubular reabsorption capacity is reduced (renal glucosuria, due to SGLT2 mutations or acquired proximal tubular dysfunction/Fanconi syndrome). SGLT2 inhibitor medications (canagliflozin, dapagliflozin, empagliflozin) pharmacologically produce glucosuria by blocking SGLT2 to lower blood glucose in type 2 diabetes. Urine dipstick for glucose is a screening test; positive results should be confirmed with fasting plasma glucose, oral glucose tolerance test, or HbA1c.
\section*{Alternative Names}
Urine Glucose; Glucosuria; Urinalysis for Glucose; Urine Sugar; Glycosuria; Renal Glucosuria
\section*{Principle}
Urine dipstick glucose oxidase/peroxidase reaction: Glucose + O2 $\rightarrow$ gluconic acid + H2O2 (glucose oxidase); H2O2 + chromogen (reduced) $\rightarrow$ oxidized chromogen (colored) + H2O (peroxidase). The dipstick is specific for glucose and does NOT react with fructose, galactose, lactose, or other reducing sugars (which are detected by the Clinitest copper reduction tablet method). False-negative dipstick results from: high concentrations of ascorbic acid (vitamin C), ketones, or levodopa, which inhibit the peroxidase reaction. False-positive: hydrogen peroxide or strong oxidizing agents contaminating the urine container \cite{rifai2018tietz}.
\section*{History}
The sweet taste of diabetic urine was known in antiquity and was confirmed by Thomas Willis in the 1670s, who gave the disease the name ``diabetes mellitus.'' Chemical copper-reduction tests for reducing sugars, including Trommer's test (1841), Fehling's test (1848), and Benedict's reagent (1908), served as the standard urine sugar methods well into the mid-20th century \cite{benedict1909reagent}. In the 1940s, Clinitest tablets brought copper-reduction testing to the bedside and clinic. Glucose-specific dipsticks based on the glucose oxidase--peroxidase reaction were introduced in the late 1950s, allowing semiquantitative urine glucose testing that remains in use today \cite{free1957simple}. The renal sodium-glucose cotransporters were identified with the cloning of SGLT1 in 1987 and SGLT2 in the early 1990s \cite{wright2011biology}, and selective SGLT2 inhibitors entered clinical use for type 2 diabetes beginning in 2013, making pharmacologic glucosuria a routine finding.
\section*{Why This Test Is Ordered}
\begin{itemize}
  \item As part of a routine urinalysis, where glucosuria may be the first clue to previously undiagnosed diabetes
  \item Opportunistic diabetes screening when blood glucose or HbA1c testing is unavailable, particularly in resource-limited settings
  \item Monitoring of glycemic control in patients who cannot measure blood glucose, although blood glucose and HbA1c are preferred
  \item Confirmation of expected glucosuria in patients on SGLT2 inhibitor therapy
  \item Evaluation of suspected renal (normoglycemic) glucosuria, including workup of Fanconi syndrome and proximal tubular dysfunction
  \item Investigation of polyuria, polydipsia, or recurrent urinary symptoms in which diabetes is in the differential diagnosis
\end{itemize}
\section*{Is This a Primary or Secondary Test}
Urine glucose is a secondary screening test and is not the primary method for diagnosing or monitoring diabetes, for which fasting plasma glucose, oral glucose tolerance testing, and HbA1c are the standards of care. Its principal roles are opportunistic screening on routine urinalysis and the evaluation of tubular (renal) glucosuria, in which blood glucose is normal.
\section*{How This Test Is Performed}
A random urine specimen is collected in a clean container, ideally by clean-catch midstream technique, and the dipstick pad is immersed and read after 30--60 seconds, either by color comparison or by an automated urine analyzer. The glucose oxidase pad is specific for glucose and does not detect other reducing sugars. Semiquantitative results are reported as trace to 4+ or in mg/dL; quantitative 24-hour urine glucose is reserved for renal glucosuria assessment. Results are available immediately at the point of care.
\section*{What Preparation Is Needed}
None; a random sample is usually acceptable, and no fasting is required. A fasting or timed postprandial sample improves interpretation by relating glucosuria to the prevailing blood glucose. Patients should avoid excessive vitamin C ingestion shortly before testing, and the sample should be tested promptly because glucose is consumed by bacteria and cells in standing urine.
\section*{Contraindications and Precautions}
There are no contraindications to urine collection. Interpretive cautions: (1) False-negative dipstick results occur with high urinary concentrations of ascorbic acid (vitamin C), ketones, or levodopa, which inhibit the peroxidase reaction. (2) False-positive results occur with peroxide contamination of the container or testing area. (3) A negative urine glucose does not exclude hyperglycemia because glucosuria appears only when the renal threshold (approximately 180 mg/dL) is exceeded. (4) Glucosuria in pregnancy is expected at lower blood glucose levels because of the physiologically reduced renal threshold, so it is a poor marker of gestational diabetes. (5) Patients on SGLT2 inhibitors have glucosuria irrespective of blood glucose, so the finding has no diagnostic significance in that setting.
\section*{Risks}
There is essentially no risk from urine collection itself; contamination of the sample is the main concern. Rarely, prolonged holding of the specimen may introduce infection or invalidate the result.
\section*{What Happens After the Test}
A positive urine glucose is confirmed with blood testing, including fasting plasma glucose, oral glucose tolerance testing, or HbA1c, before a diagnosis of diabetes is made \cite{ada2024standards}. Normoglycemic glucosuria prompts evaluation for renal tubular disorders, including a urine panel for amino acids, phosphate, uric acid, and protein (Fanconi syndrome workup) and consideration of SLC5A2 (SGLT2) gene testing. In patients taking SGLT2 inhibitors, persistent glucosuria requires no further evaluation.
\section*{Understanding Results}
\subsection*{Negative (Normal)}
No glucose detected in urine. Blood glucose below the renal threshold (typically <180 mg/dL). In a patient with diabetes, a negative urine glucose suggests good glycemic control (blood glucose <180--200 mg/dL). Patients on SGLT2 inhibitors will have glucosuria despite normal or well-controlled blood glucose.
\subsection*{Positive (Glucosuria Detected)}
Diabetes mellitus: blood glucose exceeds the renal threshold (typically >180--200 mg/dL). Uncontrolled type 1 or type 2 diabetes. Gestational diabetes: glucosuria is common in pregnancy because the renal threshold is physiologically lower (~140--160 mg/dL) due to increased GFR and reduced tubular glucose reabsorption. SGLT2 inhibitor therapy: expected glucosuria. Renal glucosuria: persistent glucosuria with normal blood glucose levels (due to SLC5A2 [SGLT2] loss-of-function mutations or acquired proximal tubular dysfunction/Fanconi syndrome, where aminoaciduria, phosphaturia, uricosuria, bicarbonaturia, and low-molecular-weight proteinuria accompany the glucosuria).
\section*{Normal Results}
Negative (dipstick): <50--100 mg/dL glucose in urine (the detection limit of the dipstick). Quantitative urine glucose: <0.5 g/24 hours (<200 mg/24 hours). Renal threshold for glucose: approximately 180 mg/dL (10 mmol/L), lower in pregnancy. Urine glucose is NOT recommended as a primary method for diabetes monitoring (blood glucose monitoring, HbA1c, and continuous glucose monitoring are the standard of care).
\section*{Follow-up Tests}
Fasting plasma glucose, oral glucose tolerance test (pregnancy), HbA1c, serum BUN/creatinine (renal function), and if renal glucosuria is suspected: spot urine for amino acids, phosphate, uric acid, protein, and bicarbonate (Fanconi syndrome workup); SLC5A2/SGLT2 gene testing \cite{santer2003molecular}.
\section*{References}
\bibliographystyle{unsrtnat}
\bibliography{references}
\clearpage
